% Page 055 — page imprimée 52
% Les Chantiers de Jeunesse — éclatement et fin des Chantiers

\subsection*{L’éclatement et la fin des Chantiers}

Il semble que le groupement 19 ait joui d’un moindre endoctrinement que certains
autres groupements et d’une direction moins pro-vichyste.

Mais dès la fin de l’année 1942 l’horizon s’assombrit pour les chantiers, car, avec le
débarquement allié en Afrique du Nord en Novembre 1942 et la capitulation de
l’armée allemande à Stalingrad le 2 Février 1943, les autorités d’occupation
demandent, puis exigent le départ des groupements pour le Service du Travail
Obligatoire. Beaucoup sont obligés de partir, parfois entraînés par leurs chefs,
drapeau et musique en tête pour travailler dans les usines d’armement allemandes. Ce
ne fut pas le cas à Meyrueis.

Marceline Causse, de Pradines, rapporte cette anecdote dans ses souvenirs 1942-
1944 :

\begin{quote}\itshape
« Frank Robert, pasteur de Meyrueis, voit à la poste 150 jeunes téléphonant à leurs familles qu’ils
partent au STO en Allemagne et les encourage à ne pas y aller. Il apprend plus tard que ces 150 sont
arrivés seulement 13 à Grenoble. Le train, traversant la plaine languedocienne, s’arrêtait, s’évadait
qui voulait ».
\end{quote}

Le groupement 19 continuera à vivoter : Il semble avoir été oublié. Il sera dissous
officiellement le 15 Juin 1944, bien longtemps après les autres groupements. Le
débarquement et les combats de la Parade dataient du 6 juin ...

\subsection*{Postface}

Après avoir rédigé ce très bref résumé de la vie des Jeunes appelés du groupement 19
à Meyrueis, quelques documents m’ont été remis qui soulignent et complètent les
lignes qui précèdent ainsi ce commentaire à la suite d’un rassemblement d’anciens
en Août 1991

\begin{quote}\itshape
« 15 Août 1941 : Il y a juste 50 ans nous descendions de la Pépinière à marche
forcée. Grand messe dans la prairie de Roquedols, face au château. Défilé dans
toutes les rues de Meyrueis avec tout le groupement 19. Rassemblement sur la Place
d’Armes. Remise des fanions aux différents groupes. Repas au pré de l’Ayrette puis
retour au camp

15 Août 1991 : Cinquante ans après, retour aux sources de notre jeunesse. Grand
messe à l’Eglise de Meyrueis et le soir procession aux flambeaux de Notre Dame du
rocher à l’Eglise de Meyrueis ... ». Jacques Cattoir 59130 Lambersart
\end{quote}

Le groupe 4 était installé pour partie à Campredon mais aussi à Valbelle et quelques
uns composèrent cette chanson qui décrit les conditions difficiles de leur vie sur le
versant nord du Calcadis à 1100m d’altitude.
